\documentclass[11pt]{article}

\usepackage[margin=1in]{geometry}
\usepackage{amsmath,amssymb,amsfonts,mathtools,amsthm}
\usepackage{bm}
\usepackage{booktabs,array,longtable,tabularx}
\usepackage{enumitem}
\usepackage{xcolor}
\usepackage{xurl}
\usepackage[hidelinks]{hyperref}

\numberwithin{equation}{section}
\setlist{nosep,leftmargin=*}
\setlength{\emergencystretch}{3em}

\newcommand{\CP}{\mathbb{CP}}
\newcommand{\CC}{\mathbb C}
\newcommand{\RR}{\mathbb R}
\newcommand{\ZZ}{\mathbb Z}
\newcommand{\cO}{\mathcal O}
\newcommand{\cH}{\mathcal H}
\newcommand{\cI}{\mathcal I}
\newcommand{\dd}{\mathrm d}
\newcommand{\ii}{\mathrm i}
\newcommand{\Tr}{\operatorname{Tr}}
\newcommand{\Sym}{\operatorname{Sym}}
\newcommand{\Span}{\operatorname{Span}}
\newcommand{\diag}{\operatorname{diag}}
\newcommand{\status}[1]{\textsf{#1}}

\definecolor{darkgreen}{RGB}{0,105,55}
\definecolor{darkred}{RGB}{150,25,25}
\definecolor{darkblue}{RGB}{20,65,140}

\newtheorem{theorem}{Theorem}[section]
\newtheorem{proposition}[theorem]{Proposition}
\newtheorem{lemma}[theorem]{Lemma}
\newtheorem{corollary}[theorem]{Corollary}
\newtheorem{definition}[theorem]{Definition}
\newtheorem{assumption}[theorem]{Assumption}
\newtheorem{gate}[theorem]{Promotion gate}
\newtheorem{remark}[theorem]{Remark}

\title{AP-E1: The Projective-Doublet Action\\
Exact Quotient Theorems, Fixed-Charge Quantization,\\
and the First-Principles Boundary of the Route-E Bridge}
\author{Route F research note}
\date{14 July 2026}

\begin{document}
\maketitle

\begin{abstract}
This note performs AP-E1 of the Another-Physics/Route-E program.  The central
result is a separation theorem.  One nonzero charge-one complex doublet with a
\emph{local} common-phase redundancy and fixed norm has the exact quotient
\(S^3/U(1)=\CP^1\); eliminating the auxiliary connection gives the
Fubini--Study sigma model.  By contrast, a physical \emph{global} common phase
cannot be removed pointwise.  At fixed global charge, a collective doublet
reduces instead to a charged rotor on \(T^*\CP^1\) in a unit Dirac-monopole
background.  Holomorphic sections \(H^0(\CP^1,\cO(k))\) arise only after a
separate first-order symplectic reduction or a controlled lowest-Landau-level
projection.

This distinction exposes a decisive Route-E blocker.  The Hopf line has Chern
number one, whereas \(T_{\CP^1}\simeq\cO(2)\).  Neither \(k=2\), the
holomorphic projection, nor the identification of its three states with three
chiral families follows from the doublet quotient.  Moreover, identifying the
prequantum level with the existing macroscopic Q-ball benchmark gives
\(k=Q/\hbar=10^6\) and \(\dim H^0(\cO(k))=1{,}000{,}001\), not three.  A
new separated-level construction is proposed: retain the macroscopic conserved
charge and attach an independent level-two projective worldline sector.  This
is mathematically consistent and falsifiable, but its level-selection and UV
origin remain explicit bridge axioms.  A 30-check arithmetic and source
regression audit verifies the quotient geometry, Chern number, Routh transform,
spin-one algebra, signed rotor spectrum, Branch-B symplectic sign, Q-ball
mismatch, and Derrick scaling without authorizing physics promotion.
\end{abstract}

\tableofcontents

\section{Scope, conventions, and final logical status}

We use spacetime signature \((-+++ )\), \(\hbar\) explicitly in the quantum
mechanics, and a single complex doublet
\begin{equation}
 Z(x)=\begin{pmatrix}z_0(x)\\z_1(x)\end{pmatrix}\in\CC^2.
 \label{eq:doublet}
\end{equation}
``One doublet'' means two complex components.  Two independent doublets form
four complex components and, under one fixed-total-norm constraint and one
common \(U(1)\) quotient, generically give \(\CP^3\), not \(\CP^1\).

Four labels are used throughout.
\begin{longtable}{>{\raggedright\arraybackslash}p{0.19\textwidth}
                  >{\raggedright\arraybackslash}p{0.72\textwidth}}
\toprule
Label & Meaning \\
\midrule
\endfirsthead
\toprule
Label & Meaning \\
\midrule
\endhead
\status{proved} & A consequence of the displayed action and hypotheses.\\
\status{no-go} & A stated set of hypotheses is incompatible with the claimed
conclusion.\\
\status{conditional} & A consistent construction after an additional
low-energy limit, projection, or sector choice.\\
\status{bridge axiom} & An identification or level-selection rule not implied
by either the projective-doublet action or Route E.\\
\bottomrule
\end{longtable}

\begin{theorem}[AP-E1 separation theorem]
The following statements hold.
\begin{enumerate}
\item A fixed-norm complex doublet with a free local common-phase action has
\(
 \{Z:Z^\dagger Z=v^2\}/U(1)\simeq\CP^1
\), and the canonical flat kinetic term induces the Fubini--Study metric.
\item If the common phase is only global and the kinetic metric is
nondegenerate, the phase is a local physical field.  Quotienting one constant
phase does not produce a pointwise \(\CP^1\) field theory.
\item Fixing a global Noether charge and reducing a collective doublet gives a
monopole-coupled \(T^*\CP^1\) rotor.  Its full Hilbert space contains infinitely
many Landau levels.
\item \(H^0(\CP^1,\cO(k))\) is obtained from first-order K\"ahler reduction or
from the lowest Landau level of that rotor.  For \(k=2\) it is the spin-one
triplet, but \(k=2\) and the projection are additional inputs.
\item The pure two-derivative \(3+1\)-dimensional \(\CP^1\) model has no
finite-size static soliton.  A four-derivative, flux, charge, wall, or gravity
mechanism is required, and the full radial/gauge theory must still pass a
stability test.
\end{enumerate}
\end{theorem}

\begin{gate}[AP-E1 non-promotion]
AP-E1 may establish the projective geometry but may not promote the Route-E
\(\cO(2)\)/three-family mechanism until a physical level-two source, a
controlled holomorphic/LLL gap, a stable background, and a chiral fluctuation
operator are all supplied.  The numerical audit therefore keeps
\path{physics_promotion_allowed=false}.
\end{gate}

\section{Minimal actions and the exact quotient}

\subsection{Auxiliary local action}

Assume a global \(SU(2)\) action on \(Z\) and a local common phase
\begin{equation}
 Z\longmapsto e^{\ii\alpha(x)}Z,
 \qquad
 a_\mu\longmapsto a_\mu+\partial_\mu\alpha,
 \qquad
 D_\mu Z=(\nabla_\mu-\ii a_\mu)Z.
 \label{eq:localgauge}
\end{equation}
The exact fixed-norm auxiliary action is
\begin{equation}
 S_{\rm aux}=-\int\dd^4x\sqrt{-g}\left[
 (D_\mu Z)^\dagger D^\mu Z+\Lambda(x)(Z^\dagger Z-v^2)\right],
 \qquad v>0.
 \label{eq:auxaction}
\end{equation}
The multiplier equation fixes \(Z^\dagger Z=v^2\), while the absence of a
Maxwell term makes \(a_\mu\) algebraic.

A renormalizable linear completion is
\begin{equation}
 S_{\rm lin}=\int\dd^4x\sqrt{-g}\left[-(D_\mu Z)^\dagger D^\mu Z
 -m^2X-\frac{\lambda}{2}X^2
 -\frac{1}{4e^2}f_{\mu\nu}f^{\mu\nu}\right],
 \quad X=Z^\dagger Z,
 \label{eq:linearaction}
\end{equation}
where \(m^2<0\), \(\lambda>0\), and
\begin{equation}
 v^2=-\frac{m^2}{\lambda}.
 \label{eq:vev}
\end{equation}
Taking the radial stiffness large freezes \(X\), while taking the vector
kinetic scale heavy makes \(a_\mu\) auxiliary at leading derivative order.
For finite \(\lambda,e\), the \(\CP^1\) theory is an effective derivative
expansion, not an exact equality of full theories.  Gauged linear models give
the standard action-level context for projective phases
\cite{Witten1993}; recent linear formulations of related Grassmannian models
illustrate, but do not prove, a four-dimensional UV completion of the present
model \cite{BykovKrivorol2024}.

\begin{remark}[What symmetry does and does not fix]
In a four-dimensional renormalizable scalar theory the canonical two-derivative
term is fixed up to normalization after the field coordinates are chosen.
Higher-dimensional EFT operators can deform the radial metric.  After the
quotient, however, every \(SU(2)\)-invariant positive metric on
\(SU(2)/U(1)\simeq\CP^1\) is proportional to the round/Fubini--Study metric:
the real two-dimensional tangent representation of the isotropy \(U(1)\) is
irreducible, so an invariant positive bilinear form is unique up to scale.
\end{remark}

\subsection{Quotient theorem}

\begin{theorem}[Hopf quotient]
\label{thm:hopfquotient}
For \(v>0\), the vacuum manifold of \eqref{eq:auxaction} modulo the local
common phase is
\begin{equation}
 \mathcal M=\{Z\in\CC^2:Z^\dagger Z=v^2\}/U(1)
 =S^3_v/U(1)\simeq\CP^1\simeq S^2.
 \label{eq:hopfquotient}
\end{equation}
\end{theorem}

\begin{proof}
The constraint surface is \(S^3_v\).  Because \(v\ne0\), the common-phase
action is free.  Map each orbit to
\begin{equation}
 P=\frac{ZZ^\dagger}{v^2},
 \qquad P^2=P,\quad P^\dagger=P,\quad \Tr P=1.
 \label{eq:projector}
\end{equation}
The projector is phase invariant.  Conversely, a rank-one Hermitian projector
selects a complex line in \(\CC^2\), and every unit vector spanning that line
differs only by a phase.  Rank-one complex lines are precisely
\(\mathbb P(\CC^2)=\CP^1\).  Writing
\begin{equation}
 P=\frac12(I+n^a\sigma_a),
 \qquad n^a=\frac{Z^\dagger\sigma^a Z}{v^2},
 \label{eq:hopfmap}
\end{equation}
the Pauli completeness identity
\begin{equation}
 (\sigma^a)_{ij}(\sigma^a)_{k\ell}
 =2\delta_{i\ell}\delta_{jk}-\delta_{ij}\delta_{k\ell}
\end{equation}
gives \(n^an^a=1\), identifying the quotient with \(S^2\).
\end{proof}

This is also the elementary Marsden--Weinstein quotient
\(\CC^2//U(1)\) at a nonzero moment-map value
\cite{MarsdenWeinstein1974}.  It proves the quotient after the field content,
charge, and nonzero level are stated; it does not derive those inputs.

\section{Flat-doublet polar decomposition and induced geometry}

\subsection{One exact identity}

On the patch \(z_0\ne0\), set
\begin{equation}
 N=1+|w|^2,
 \qquad
 s_N(w)=\frac{(1,w)^T}{\sqrt N},
 \qquad
 Z=\rho e^{\ii\chi}s_N(w).
 \label{eq:northsection}
\end{equation}
Define the real Hopf connection and Fubini--Study quadratic form by
\begin{align}
 \mathcal A&=-\ii s_N^\dagger\dd s_N
 =\frac{\ii(w\dd\bar w-\bar w\dd w)}{2N},
 \label{eq:hopfconnection}\\
 \dd\ell_{\rm FS}^2&=
 \dd s_N^\dagger(I-s_Ns_N^\dagger)\dd s_N
 =\frac{\dd\bar w\,\dd w}{N^2}.
 \label{eq:fsmetric}
\end{align}

\begin{theorem}[Exact polar decomposition]
The flat Hermitian metric on \(\CC^2\setminus\{0\}\) obeys
\begin{equation}
 \dd Z^\dagger\dd Z
 =(\dd\rho)^2+\rho^2\left[(\dd\chi+\mathcal A)^2
 +\dd\ell_{\rm FS}^2\right].
 \label{eq:polaridentity}
\end{equation}
\end{theorem}

\begin{proof}
Since \(s_N^\dagger s_N=1\), differentiation gives
\(s_N^\dagger\dd s_N\) purely imaginary and
\(s_N^\dagger\dd s_N=\ii\mathcal A\).  Decompose
\begin{equation}
 \dd s_N=s_N(s_N^\dagger\dd s_N)
 +(I-s_Ns_N^\dagger)\dd s_N.
\end{equation}
The two terms are orthogonal, so
\begin{equation}
 \dd s_N^\dagger\dd s_N
 =\mathcal A^2+\dd\ell_{\rm FS}^2.
\end{equation}
Substitution in
\(
 \dd Z=e^{\ii\chi}[s_N\dd\rho+\rho(\ii s_N\dd\chi+\dd s_N)]
\)
gives \eqref{eq:polaridentity}; all radial cross terms cancel.
\end{proof}

Thus the connection and the metric are not appended geometric decorations:
they are respectively the vertical and horizontal parts of the flat doublet
metric.  In spherical coordinates
\begin{equation}
 w=e^{\ii\phi}\tan\frac\theta2,
\end{equation}
one obtains the convention-fixed formulas
\begin{align}
 \mathcal A_N&=\frac{1-\cos\theta}{2}\dd\phi,
 \label{eq:AN}\\
 \mathcal F=\dd\mathcal A_N
 &=\frac12\sin\theta\,\dd\theta\wedge\dd\phi,
 \label{eq:curvature}\\
 \dd\ell_{\rm FS}^2
 &=\frac14(\dd\theta^2+\sin^2\theta\,\dd\phi^2)
 =\frac14\dd\bm n\cdot\dd\bm n.
 \label{eq:fssphere}
\end{align}
Our real line element therefore has sphere radius \(1/2\), Gaussian curvature
four, and scalar curvature eight.  Other common complex-coordinate
conventions place a factor of two in \(\dd s^2=2g_{w\bar w}\dd w\dd\bar w\);
all action, flux, and spectrum formulas here use \eqref{eq:fsmetric}.

\subsection{Patch transition and Chern number}

On \(z_1\ne0\), let \(\eta=1/w\) and
\begin{equation}
 s_S(\eta)=\frac{(\eta,1)^T}{\sqrt{1+|\eta|^2}}.
\end{equation}
On the overlap, with \(w=|w|e^{\ii\phi}\),
\begin{equation}
 s_S=e^{-\ii\phi}s_N,
 \qquad
 \mathcal A_S=\mathcal A_N-\dd\phi.
 \label{eq:patchtransition}
\end{equation}
There is no global smooth Hopf section, but \(\mathcal F\) is global and
\begin{equation}
 \frac{1}{2\pi}\int_{\CP^1}\mathcal F
 =\frac{1}{2\pi}\int_0^{2\pi}\dd\phi\int_0^\pi
 \frac12\sin\theta\,\dd\theta=1.
 \label{eq:chernone}
\end{equation}
This is the unit Hopf bundle.  The sign of the line-bundle name depends on
whether one associates charge \(+1\) or \(-1\) to the principal bundle; we fix
the convention that the covariant derivative \(D=\dd-\ii k\mathcal A\) has
\(c_1=k\) and call it \(\cO(k)\).

\subsection{Eliminating the auxiliary field}

Varying \eqref{eq:auxaction} with respect to \(a_\mu\) gives
\begin{equation}
 a_\mu^\star=-\frac{\ii}{2Z^\dagger Z}
 \left[Z^\dagger\partial_\mu Z
 -(\partial_\mu Z^\dagger)Z\right].
 \label{eq:astargeneral}
\end{equation}
At fixed norm this becomes
\begin{equation}
 a_\mu^\star=-\frac{\ii}{v^2}Z^\dagger\partial_\mu Z
 =\partial_\mu\chi+\mathcal A_\mu.
 \label{eq:astar}
\end{equation}
Substitution projects orthogonally to the gauge orbit:
\begin{align}
 (D_\mu Z)^\dagger D^\mu Z\big|_{a^\star}
 &=\partial_\mu Z^\dagger
 \left(I-\frac{ZZ^\dagger}{v^2}\right)\partial^\mu Z
 \notag\\
 &=v^2\frac{\partial_\mu\bar w\,\partial^\mu w}
 {(1+|w|^2)^2}
 =\frac{v^2}{4}\partial_\mu\bm n\cdot\partial^\mu\bm n.
 \label{eq:horizontalaction}
\end{align}
Hence
\begin{equation}
 S_{\CP^1}=-v^2\int\dd^4x\sqrt{-g}\,
 \frac{\partial_\mu\bar w\,\partial^\mu w}{(1+|w|^2)^2}.
 \label{eq:cp1action}
\end{equation}
Its affine equation of motion is
\begin{equation}
 \Box w-\frac{2\bar w}{1+|w|^2}
 (\partial_\mu w)(\partial^\mu w)=0.
 \label{eq:cpeom}
\end{equation}
This classical projective construction is standard in the
\(\CP^{N-1}\) sigma-model literature \cite{DAdda1978}.

\section{How far the derivation is genuinely first-principles}

There are two independent routes to the same metric, and their common boundary
should be explicit.

\subsection{Dynamical-field route}

Assume:
\begin{enumerate}
\item a local, Lorentz-invariant, positive two-derivative theory;
\item exactly one complex order-parameter doublet with a global \(SU(2)\)
basis symmetry;
\item common phase is a local redundancy;
\item the radial mode is restricted to a nonzero level.
\end{enumerate}
Then Theorem~\ref{thm:hopfquotient} and \eqref{eq:horizontalaction} force
\(\CP^1\) and the Fubini--Study kinetic term up to the stiffness \(v^2\).
This is an action-level derivation, but the four assumptions are microscopic
inputs.

\subsection{Ray/fidelity route}

Alternatively, suppose each point carries a normalized binary pure state
\(|u\rangle\in\CC^2\), physical probabilities identify
\(|u\rangle\sim e^{\ii\alpha}|u\rangle\), and nearby states are compared by
fidelity.  Expanding
\begin{equation}
 1-|\langle u|u+\dd u\rangle|^2
\end{equation}
to quadratic order gives
\begin{equation}
 \dd s_{\rm fidelity}^2
 =\langle\dd u|\dd u\rangle
 -|\langle u|\dd u\rangle|^2,
 \label{eq:fidelity}
\end{equation}
which is exactly \eqref{eq:fsmetric}.  Pulling this information metric back to
spacetime yields \eqref{eq:cp1action}.  This route explains why phase is a ray
redundancy, but it still assumes a binary pure state and does not show that the
ray is a fundamental spacetime field.

\begin{corollary}[First-principles boundary]
Locality, Lorentz symmetry, positivity, and \(SU(2)\) invariance determine the
form of the low-energy action only \emph{after} one complex doublet, ray
equivalence, and a nonzero radial level are selected.  They do not derive the
number two, the scale \(v\), or the microscopic origin of the order parameter.
\end{corollary}

\section{Global-only no-go theorem}

Suppose instead that the common phase is a physical global symmetry.  At fixed
norm, \eqref{eq:polaridentity} gives the exact local kinetic decomposition
\begin{equation}
 \partial_\mu Z^\dagger\partial^\mu Z
 =v^2\left[
 (\partial_\mu\chi+\mathcal A_\mu)^2
 +\frac{\partial_\mu\bar w\,\partial^\mu w}{N^2}
 \right].
 \label{eq:globalkinetic}
\end{equation}

\begin{theorem}[Global \(U(1)\) cannot give a pointwise projective target]
Assume the common phase is global, the kinetic metric is positive and
nondegenerate, only one constant \(U(1)\) orbit is quotiented, and no further
constraint or gauge redundancy is introduced.  Then the local target is not
\(\CP^1\).
\end{theorem}

\begin{proof}
Equation \eqref{eq:globalkinetic} contains the independent local field
\(\chi(x)\).  Quotienting \(\chi\mapsto\chi+\alpha_0\) removes one constant
zero mode, not arbitrary functions \(\chi(x)\).  Equivalently, the fixed-norm
global theory has target \(S^3\), hence three local real coordinates, whereas
\(\CP^1\) has two.  The Euler--Lagrange equation
\begin{equation}
 \Box\chi=-\nabla_\mu\mathcal A^\mu
 \label{eq:chieom}
\end{equation}
shows that eliminating \(\chi\) requires an inverse wave operator.  Formally,
it produces a nonlocal transverse term
\begin{equation}
 \mathcal A_\mu\left(g^{\mu\nu}
 -\frac{\nabla^\mu\nabla^\nu}{\Box}\right)\mathcal A_\nu,
 \label{eq:nonlocalterm}
\end{equation}
not the pure local Fubini--Study action.
\end{proof}

This is the first no-free-lunch statement: the same common phase cannot be
simultaneously discarded as a local gauge redundancy and retained as the
physical global charge that stabilizes an ordinary Q-ball.

\section{Fixed global charge: exact collective Routh reduction}

\subsection{Doublet Q-ball ansatz}

Consider a global \(U(2)\)-invariant scalar theory in flat spacetime,
\begin{equation}
 S_{\rm gl}=\int\dd^4x\left[-\partial_\mu\Phi^\dagger\partial^\mu\Phi
 -U(\Phi^\dagger\Phi)\right].
 \label{eq:globalaction}
\end{equation}
For the spatially coherent orientation ansatz
\begin{equation}
 \Phi(\bm x,t)=f(\bm x)e^{\ii\chi(t)}s_N(w(t)),
 \label{eq:collectiveansatz}
\end{equation}
define
\begin{equation}
 \cI[f]=2\int\dd^3x\,f^2,
 \qquad
 E_0[f]=\int\dd^3x\left[(\bm\nabla f)^2+U(f^2)\right].
 \label{eq:inertia}
\end{equation}
Direct substitution, without using field equations, gives
\begin{equation}
 L_{\rm coll}=\frac{\cI}{2}\left[
 (\dot\chi+\mathcal A_i\dot q^i)^2
 +g^{\rm FS}_{ij}\dot q^i\dot q^j\right]-E_0,
 \label{eq:lcoll}
\end{equation}
where \(q^i\) are real coordinates on \(\CP^1\).  The global charge is
\begin{equation}
 Q=-\ii\int\dd^3x
 (\Phi^\dagger\dot\Phi-\dot\Phi^\dagger\Phi)
 =\cI(\dot\chi+\mathcal A_i\dot q^i).
 \label{eq:Qglobal}
\end{equation}
Consequently the fixed-profile energy, before minimizing the orientation
velocity, is
\begin{equation}
 E[f,q,\dot q]=E_0[f]+\frac{Q^2}{2\cI[f]}
 +\frac{\cI[f]}{2}g^{\rm FS}_{ij}\dot q^i\dot q^j.
 \label{eq:fixedQenergyfull}
\end{equation}
Thus the orientation-ground-state value (equivalently, the minimum over
\(\dot q\)) is
\begin{equation}
 E_Q^{\rm min}[f]=\min_{\dot q}E[f,q,\dot q]
 =E_0[f]+\frac{Q^2}{2\cI[f]}
 =E_0[f]+\frac{Q^2}{4\int f^2}.
 \label{eq:fixedQenergy}
\end{equation}
This is compatible with the standard Q-ball variational principle
\cite{Coleman1985APE1,HeeckSokhashvili2023,EspinosaHeeckSokhashvili2023}.
It is exact after insertion of \eqref{eq:collectiveansatz}, but the ansatz is
not yet proved to be a consistent truncation: spatially varying orientation,
radial radiation, and nonspherical modes have been omitted.  Fixed-charge EFT
is independently testable in other theories, as recent lattice work confirms
for the three-dimensional critical \(O(2)\) model \cite{CuomoEtAl2024}; that
evidence does not determine the coefficients of the present model.

\subsection{Routh theorem}

\begin{theorem}[Fixed-charge projective rotor]
For
\begin{equation}
 L=\frac{\cI}{2}\left[(\dot\chi+\mathcal A_i\dot q^i)^2
 +g_{ij}\dot q^i\dot q^j\right]-V(q),
\end{equation}
the Routhian at fixed \(Q\) and the reduced Hamiltonian are
\begin{align}
 R_Q&=L-Q\dot\chi
 =\frac{\cI}{2}g_{ij}\dot q^i\dot q^j
 +Q\mathcal A_i\dot q^i-V-\frac{Q^2}{2\cI},
 \label{eq:routh}\\
 H_Q&=\frac{1}{2\cI}g^{ij}
 (p_i-Q\mathcal A_i)(p_j-Q\mathcal A_j)
 +V+\frac{Q^2}{2\cI}.
 \label{eq:hamiltonian}
\end{align}
\end{theorem}

\begin{proof}
Solve \eqref{eq:Qglobal} for
\begin{equation}
 \dot\chi=\frac{Q}{\cI}-\mathcal A_i\dot q^i
\end{equation}
and substitute into \(L-Q\dot\chi\), giving \eqref{eq:routh}.
Then
\begin{equation}
 p_i=\frac{\partial R_Q}{\partial\dot q^i}
 =\cI g_{ij}\dot q^j+Q\mathcal A_i,
\end{equation}
and the Legendre transform yields \eqref{eq:hamiltonian}.
\end{proof}

The reduced classical phase space is the magnetic cotangent bundle
\(T^*\CP^1\), not merely \(\CP^1\).  For compact \(\chi\sim\chi+2\pi\),
single-valued wavefunctions give
\begin{equation}
 Q=\hbar k,
 \qquad k\in\ZZ.
 \label{eq:chargequantization}
\end{equation}
Under \(s\mapsto e^{\ii\beta}s\), one has
\(\mathcal A\mapsto\mathcal A+\dd\beta\) and
\(\chi\mapsto\chi-\beta\).  The reduced wavefunction is therefore a section
with covariant derivative
\begin{equation}
 D=\dd-\ii k\mathcal A,
 \qquad D^2=-\ii k\mathcal F,
 \qquad \frac{\ii}{2\pi}\int D^2=k.
 \label{eq:lineconnection}
\end{equation}

\section{Quantization, line bundles, and the Route-E obstruction}

\subsection{First-order symplectic reduction}

Let \(b=(b_0,b_1)^T\in\CC^2\) have symplectic form
\begin{equation}
 \Omega=\ii\hbar\sum_{A=0}^1\dd b_A\wedge\dd\bar b_A,
 \qquad \mu=\hbar b^\dagger b.
 \label{eq:symplecticC2}
\end{equation}
For \(k>0\), at \(\mu=k\hbar\), write
\begin{equation}
 b=\sqrt{k}\,e^{\ii\chi}s_N(w).
\end{equation}
The canonical one-form restricts, up to an exact term, to
\(k\hbar(\dd\chi+\mathcal A)\), hence
\begin{equation}
 \mu^{-1}(k\hbar)/U(1)=\CP^1\qquad(k>0),
 \qquad
 \Omega_{\rm red}=k\hbar\mathcal F.
 \label{eq:MWreduction}
\end{equation}
The zero level is exceptional:
\begin{equation}
 \mu^{-1}(0)=\{0\},\qquad \mu^{-1}(0)/U(1)=\{\mathrm{pt}\};
\end{equation}
it supplies the one-dimensional trivial representation, not a geometric
\(\CP^1\) orbit.  Negative charge is obtained from the conjugate doublet (or
the opposite symplectic orientation), and uses the anti-holomorphic
polarization.  K\"ahler quantization for \(k>0\), with \(k=0\) treated by the
preceding trivial orbit, gives
\begin{equation}
 \cH_k=H^0(\CP^1,\cO(k))
 \simeq\Sym^k(\CC^2)^*,
 \qquad \dim_\CC\cH_k=k+1.
 \label{eq:borelweil}
\end{equation}
This is the clean finite-dimensional route.  In Schwinger-boson language,
\begin{equation}
 N=a_A^\dagger a_A=k,
 \qquad J_i=\frac\hbar2a^\dagger\sigma_i a,
 \qquad J^2=\frac{\hbar^2}{4}N(N+2).
 \label{eq:schwinger}
\end{equation}
For \(k=2\),
\begin{equation}
 \cH_2=\Span\{|2,0\rangle,|1,1\rangle,|0,2\rangle\},
 \quad j=1,\quad \dim\cH_2=3,\quad J^2=2\hbar^2.
 \label{eq:k2triplet}
\end{equation}
The weights of \(2J_3/\hbar\) are \((+2,0,-2)\), exactly the Route-E
weight pattern.  This is a representation-theoretic match, not yet a physical
family derivation.

\subsection{Second-order rotor and the LLL gate}

The second-order Hamiltonian \eqref{eq:hamiltonian} is different: it has all
monopole harmonics \cite{WuYang1976,BykovSmilga2023}.  Since
\eqref{eq:fssphere} is a sphere of radius \(R=1/2\), define
\(\kappa=|k|\).  The magnitude of the monopole strength is
\(S=\kappa/2\), and
\begin{equation}
 j=\frac{\kappa}{2}+n,
 \qquad n=0,1,2,\ldots.
 \label{eq:jn}
\end{equation}
The degeneracy and kinetic energy are
\begin{align}
 d_n&=2j+1=\kappa+2n+1,
 \label{eq:degeneracy}\\
 E_n^{\rm kin}
 &=\frac{\hbar^2}{2\cI R^2}
 \left[j(j+1)-\frac{\kappa^2}{4}\right]
 \notag\\
 &=\frac{2\hbar^2}{\cI}
 \left[n(n+\kappa+1)+\frac{\kappa}{2}\right].
 \label{eq:rotorspectrum}
\end{align}
Only the \(n=0\) lowest Landau level is isomorphic, for positive flux, to
\(H^0(\CP^1,\cO(k))\).  Negative flux uses the conjugate anti-holomorphic
polarization and has the same spectrum with \(\kappa=|k|\).  The gap is
\begin{equation}
 \Delta_{\rm LLL}=E_1-E_0
 =\frac{2\hbar^2(\kappa+2)}{\cI}.
 \label{eq:LLLgap}
\end{equation}
For \(k=2\), \(d_0=3\) and
\(\Delta_{\rm LLL}=8\hbar^2/\cI\).  A finite-dimensional family carrier is
therefore justified only when every relevant mixing, temperature, curvature,
and portal scale is parametrically below this gap.  Otherwise higher
\(d_n=5,7,\ldots\) multiplets participate.

\subsection{Why the projective line does not imply a level-two bundle}

The Euler sequence on \(\CP^1\) is
\begin{equation}
 0\longrightarrow\cO
 \longrightarrow\cO(1)^{\oplus2}
 \longrightarrow T_{\CP^1}\longrightarrow0,
 \label{eq:eulersequence}
\end{equation}
so taking first Chern classes gives
\begin{equation}
 T_{\CP^1}\simeq\cO(2),
 \qquad c_1(T_{\CP^1})=2.
 \label{eq:tangento2}
\end{equation}
But the minimal Hopf construction has \(c_1=1\) by \eqref{eq:chernone}.
Therefore
\begin{equation}
 \boxed{\text{projective doublet}\ \Longrightarrow\ \CP^1
 \quad\text{but not}\quad
 \text{physical prequantum bundle}=T_{\CP^1}=\cO(2).}
 \label{eq:o2obstruction}
\end{equation}
One may impose that the physical fluctuation is tangent-valued, select a
charge-two sector, or introduce an independent level-two Wess--Zumino term.
Each is a new physical input.  Coherent-state/Borel--Weil constructions explain
the representation once the level is selected \cite{Perelomov1972}; they do
not select the level.

\subsection{Macroscopic-charge contradiction}

The existing bridge benchmark uses
\begin{equation}
 Q=10^6,\qquad f_0^2=0.5\,\mathrm{GeV}^2,\qquad
 R=64.9712654357\,\mathrm{GeV}^{-1}.
\end{equation}
In the step-profile collective ansatz,
\begin{equation}
 \cI=2f_0^2\frac{4\pi R^3}{3}
 =1.1488215825\times10^6\,\mathrm{GeV}^{-1},
 \qquad
 \frac{Q}{\cI}=0.870457184321\,\mathrm{GeV},
 \label{eq:benchmarkinertia}
\end{equation}
reproducing the stored angular frequency.  If this same \(Q\) is used in
\eqref{eq:chargequantization}, then
\begin{equation}
 k=10^6,
 \qquad
 \dim H^0(\CP^1,\cO(k))=1{,}000{,}001.
 \label{eq:macrocontradiction}
\end{equation}
It cannot simultaneously be the \(k=2\) triplet.  This is not a small
correction: it separates two distinct physical sectors.

\section{Static stability and the limits of topology}

For a unit vector \(\bm n:\RR^3\to S^2\), the two-derivative energy is
\begin{equation}
 E_2=\frac{v^2}{4}\int\dd^3x\,(\partial_i\bm n)^2.
 \label{eq:E2}
\end{equation}
Use the size dilation \(\bm n_\lambda(\bm x)=\bm n(\bm x/\lambda)\).  In three
spatial dimensions,
\begin{equation}
 E_2(\lambda)=\lambda E_2,
 \qquad E_0(\lambda)=\lambda^3E_0.
 \label{eq:derricktwo}
\end{equation}
Thus a positive \(E_2+E_0\) has
\begin{equation}
 \left.\frac{\dd E}{\dd\lambda}\right|_{\lambda=1}
 =E_2+3E_0>0,
 \label{eq:derrickslope}
\end{equation}
and shrinks; this is Derrick's obstruction \cite{Derrick1964}.

If finite gauge stiffness generates, at leading order,
\begin{align}
 \mathcal F_{\mu\nu}
 &=\frac12\bm n\cdot
 (\partial_\mu\bm n\times\partial_\nu\bm n),
 \label{eq:Fhopf}\\
 \mathcal L_4
 &=-\frac{1}{4e^2}\mathcal F_{\mu\nu}\mathcal F^{\mu\nu}
 =-\frac{1}{16e^2}
 [\bm n\cdot(\partial_\mu\bm n\times\partial_\nu\bm n)]^2,
 \label{eq:skyrme}
\end{align}
then \(E_4(\lambda)=\lambda^{-1}E_4\), and
\begin{equation}
 E_2(\lambda)+E_4(\lambda)
 =\lambda E_2+\lambda^{-1}E_4
\end{equation}
has a stationary minimum at
\begin{equation}
 \lambda_\star=\sqrt{E_4/E_2}.
 \label{eq:balance}
\end{equation}
Finite-energy boundary conditions compactify space to \(S^3\), so
\begin{equation}
 \pi_3(S^2)=\ZZ,
 \qquad
 Q_H=\frac{1}{4\pi^2}\int_{S^3}\mathcal A\wedge\mathcal F.
 \label{eq:hopfcharge}
\end{equation}
The four-derivative model supports the familiar Hopf-energy mechanism and
bounds \cite{VakulenkoKapitanskii1979,FaddeevNiemi1997}.  Two-component
charged condensates provide an action-level route to similar low-energy
structures \cite{BabaevFaddeevNiemi2002}; recent preprint work studies local
stability of large twisted configurations in two-flavor scalar QED
\cite{SheuShifman2026}.  These results motivate a numerical AP-E5 gate but do
not constitute a global existence theorem for the present parameters.

Three caveats remain.
\begin{enumerate}
\item \(\pi_0(\CP^1)=0\), so the two poles are not topologically distinct
vacua.  An anisotropy such as
\(V_{\rm aniso}=\Lambda_b^4(1-n_3^2)\) creates two minima but explicitly
breaks \(SU(2)\).
\item If the radial mode can reach \(Z=0\), the projective coordinate is
undefined and Hopf charge can unwind through a core.
\item Portal forces can create an upper stability charge; conserved \(U(1)\)
alone does not guarantee arbitrarily large stable Q-balls
\cite{HamadaEtAl2024APE1}.
\end{enumerate}
Exact Hopfion solutions in an inhomogeneous magnetic Hamiltonian further show
that a length-setting deformation can evade scaling, but that condensed-matter
Hamiltonian is not a Lorentz-invariant completion of AP-E1
\cite{BalakrishnanEtAl2023}.

\section{Three completion branches}

\subsection{Branch A: physical charge two}

Set the physical compact charge sector to \(k=2\).  Equations
\eqref{eq:borelweil}--\eqref{eq:k2triplet} then give the triplet with no extra
level field.  This is the mathematically shortest route, but a two-quantum
object is not in the macroscopic thin-wall regime of the current
\(Q=10^6\) benchmark.  It needs a genuine few-body/nonperturbative calculation
rather than the classical Q-ball approximation.  Non-Abelian Q-ball analyses
show that internal charge sectors can be nontrivial
\cite{SafianColemanAxenides1988}, but they do not select this Route-E sector.

\subsection{Branch B: separated macroscopic and projective charges}

Retain a macroscopic Q-ball field \(\Phi_Q\) with global charge \(Q\), but
introduce an independent worldline doublet \(b(t)\) with local
\(U(1)_b\).  A minimal first-order action is
\begin{align}
 S_{\rm sep}&=S_Q[\Phi_Q]+\int\dd t\,L_b+S_{\rm portal},
 \label{eq:separatedaction}\\
 L_b&=-\frac{\ii\hbar}{2}
 \left[b^\dagger D_tb-(D_tb)^\dagger b\right]+\hbar k c_t,
 \qquad D_tb=(\partial_t-\ii c_t)b.
 \label{eq:firstorderb}
\end{align}
Under
\begin{equation}
 b\mapsto e^{\ii\alpha(t)}b,
 \qquad c_t\mapsto c_t+\dot\alpha,
\end{equation}
the action changes only by the endpoint term
\(+\hbar k[\alpha(t_f)-\alpha(t_i)]\).  The \(c_t\) equation imposes
\begin{equation}
 b^\dagger b=k.
\end{equation}
For closed histories, invariance of \(e^{\ii S/\hbar}\) requires
\(k\in\ZZ\).  Taking \(k=2\) gives \(\cO(2)\) independently of the
macroscopic \(Q\).  Gauge-invariant portals may depend on
\begin{equation}
 n_b^a=\frac{b^\dagger\sigma^a b}{b^\dagger b}
\end{equation}
but must preserve the separate physical \(U(1)_Q\) if the Q-ball is to remain
stable.

This branch resolves the algebraic contradiction \eqref{eq:macrocontradiction}
and is the recommended AP-E1 continuation.  It does \emph{not} explain why the
worldline level equals two or why this spin-one space carries chiral families.
A pair-only microscopic constituent rule could make two the smallest nonzero
allowed occupation, but that requires an explicit charge-two or discrete
selection rule; it may not be asserted from \(b\mapsto-b\) alone.

\subsection{Branch C: tangent self-carriage}

Postulate that the physical zero mode is a tangent deformation of the
projective background.  Then \eqref{eq:tangento2} forces \(\cO(2)\) and its
three holomorphic sections.  This is geometrically economical but moves the
missing physics into the assertion ``family mode equals tangent mode.''  AP-E4
must derive that statement from a Dirac/fluctuation operator and exclude
additional modes.

\begin{longtable}{>{\raggedright\arraybackslash}p{0.14\textwidth}
                  >{\raggedright\arraybackslash}p{0.28\textwidth}
                  >{\raggedright\arraybackslash}p{0.48\textwidth}}
\toprule
Branch & Exact gain & Remaining decisive test \\
\midrule
\endfirsthead
\toprule
Branch & Exact gain & Remaining decisive test \\
\midrule
\endhead
A & One charge sector directly gives \(H^0(\cO(2))\). & Replace the
macroscopic Q-ball approximation by a controlled two-quantum calculation.\\
B & Separates Q-ball stability charge from projective Chern level. & Derive
the level-two worldline sector and show all portals remain below the LLL and
charge-leakage gaps.\\
C & Euler sequence forces tangent bundle \(\cO(2)\). & Derive a chiral
tangent-valued zero mode with no unwanted partners.\\
\bottomrule
\end{longtable}

\section{Reproducible numerical and algebraic audit}

The script
\path{route_f/code/verify_ap_e1_projective_doublet.py} writes
\path{route_f/output/ap_e1_projective_doublet.json} and
\path{route_f/output/ap_e1_projective_doublet.md}.  Its current result is
\begin{equation}
 \boxed{30/30\ \text{checks pass},\qquad
 \texttt{physics\_promotion\_allowed=false}.}
 \label{eq:auditstatus}
\end{equation}
The checks include:
\begin{enumerate}
\item projector idempotence, Hermiticity, phase invariance, and the Hopf
\(S^2\) reconstruction;
\item exact vertical/horizontal metric decomposition, north/south chart
agreement, and gauge covariance;
\item midpoint integration
\begin{equation}
 \int\mathcal F=6.283185468670607,
 \qquad \frac{1}{2\pi}\int\mathcal F=1.000000025702;
\end{equation}
\item direct versus closed-form Routh and Legendre transforms;
\item the fixed-charge orientation-energy minimum, positive-level reduction
domain, and charge-conjugation invariance of the signed spectrum;
\item \(\dim\Sym^k\CC^2=k+1\) and its weight lists for \(0\le k\le6\);
\item the \(k=2\) commutators and Casimir \(J^2=2I_3\) to machine precision;
\item the Branch-B constraint and positive endpoint phase from its corrected
first-order symplectic sign, plus a critical-source token/hash regression;
\item the first four monopole-rotor levels and
\(\Delta_{\rm LLL}\cI/\hbar^2=8\);
\item reconstruction of the stored Q-ball frequency and the
\(1{,}000{,}001\)-versus-three contradiction;
\item the independent \(E_2,E_4,E_0\) Derrick scalings.
\end{enumerate}
These are arithmetic and source regressions, not a proof checker for every TeX
sentence.  They test convention factors and known failure modes; they cannot
establish a microscopic level-selection rule.

\section{Claim ledger and ordered next steps}

\begin{longtable}{>{\raggedright\arraybackslash}p{0.29\textwidth}
                  >{\raggedright\arraybackslash}p{0.18\textwidth}
                  >{\raggedright\arraybackslash}p{0.43\textwidth}}
\toprule
Claim & AP-E1 status & Reason \\
\midrule
\endfirsthead
\toprule
Claim & AP-E1 status & Reason \\
\midrule
\endhead
One local fixed-norm doublet gives \(\CP^1\) & \status{proved} & Free Hopf
quotient and explicit projector bijection.\\
Canonical kinetic term gives FS geometry & \status{proved} & Exact flat-metric
polar decomposition and auxiliary-field elimination.\\
Global-only phase gives pointwise \(\CP^1\) & \status{refuted} & Physical
fiber mode remains; formal elimination is nonlocal.\\
Fixed \(Q\) gives a monopole rotor & \status{proved conditionally} & Exact in
the coherent-orientation collective ansatz; full field-space truncation remains
to be justified.\\
Rotor Hilbert space equals three states & \status{refuted without LLL gate} &
Higher Landau levels give infinitely many states.\\
Doublet quotient forces \(\cO(2)\) & \status{refuted} & Minimal Hopf Chern
number is one; tangent Chern number is two.\\
\(k=2\) gives a triplet & \status{proved} & Borel--Weil/Schwinger algebra after
the level is selected.\\
Q-ball benchmark supplies \(k=2\) & \status{refuted} & Existing
\(Q=10^6\) would give dimension \(1{,}000{,}001\).\\
Triplet equals three chiral families & \status{bridge axiom} & Requires the
AP-E4 chiral operator, index, and spectral audit.\\
Pure \(\CP^1\) bubble is stable & \status{refuted} & Derrick scaling; topology
alone does not fix size.\\
Separated level-two sector is consistent & \status{conditional} & Action
\eqref{eq:separatedaction} is gauge-consistent, but its UV origin is open.\\
\bottomrule
\end{longtable}

The recommended execution order is:
\begin{enumerate}
\item \textbf{AP-E2 regression:} retain the already proved
discriminant--Killing/contact identity as an exact invariant check.
\item \textbf{AP-E3 level derivation:} implement Branch B and search for an
anomaly, pair-selection, or compactification mechanism that fixes \(k=2\).
Compute the actual \(\Delta_{\rm LLL}\) and reject the reduction if any portal
or background scale is comparable to it.
\item \textbf{AP-E4 chirality:} construct the full fluctuation/Dirac operator
on the candidate background; compute its index, kernel, partner modes, and
spectral gap.  A three-dimensional bosonic spin multiplet is not yet three
chiral fermion families.
\item \textbf{AP-E5 stability:} solve the coupled radial, gauge, and orientation
boundary-value problem at finite \(e,\lambda\).  Measure the full Hessian,
charge-emission thresholds, and the path through \(Z=0\).
\item \textbf{AP-E6 onward:} only after AP-E3--5 pass may the projective sector
be inserted into the Route-E messenger and phenomenology pipeline.
\end{enumerate}

\section{Conclusion}

AP-E1 is closed at the geometry level and open at the physical-selection
level.  The strongest unconditional statement available from the displayed
assumptions is
\begin{equation}
 \boxed{
 \text{one local fixed-norm complex doublet}
 \Longrightarrow \CP^1\ \text{with Fubini--Study dynamics}.}
\end{equation}
The strongest no-go statement is
\begin{equation}
 \boxed{
 \text{one global Q-ball phase}
 \not\Longrightarrow
 \text{pointwise }\CP^1\not\Longrightarrow\cO(2)
 \not\Longrightarrow\text{three families}.}
\end{equation}
The proposed separated-level action is the most coherent next hypothesis:
it preserves the macroscopic stability charge while making the projective
Chern level independently quantized.  Its value \(k=2\), LLL isolation,
chiral interpretation, and soliton stability are now sharply stated
falsifiable gates rather than hidden assumptions.

\bibliographystyle{unsrt}
\bibliography{route_f/tex/ap_e1_projective_doublet_action}

\end{document}
